\subsubsection{Subproceso 2: Selección y Registro de Cursos}



\begin{itemize}
    \item \textbf{Estudiante acreditado y habilitado con turno activo:} Condición emitida por el Subproceso 1 que certifica la ausencia de bloqueos administrativos, financieros o académicos y otorga al estudiante una franja horaria de acceso al portal. Sin este estado, la Intranet DIRCE impide el ingreso al módulo de selección de cursos.

    \item \textbf{Oferta académica parametrizada de la Facultad:} Conjunto de asignaturas, secciones, horarios, docentes asignados y cupos disponibles ingresados al SIGA-ORCE por los Departamentos Académicos y la Comisión de Matrícula previo al inicio del proceso.

    \item \textbf{Tabla de límites de creditaje por promedio:} Parámetro normativo (RR-0570-2022) que define la cantidad máxima de créditos que el estudiante puede registrar en función de su promedio ponderado de los dos últimos semestres académicos.

    \item \textbf{Ficha de matrícula preconfigurada (Plan de Estudios individual):} Registro en el sistema del historial académico del estudiante, incluyendo las asignaturas ya aprobadas, los prerrequisitos cumplidos y las que aún quedan pendientes, lo que determina cuáles asignaturas son visibles y seleccionables en la interfaz.

    \item \textbf{Registro de cursos obligatorios de arrastre y Grupo Cero:} Nómina de asignaturas desaprobadas del ciclo más atrasado y de cursos en condición de riesgo académico que deben priorizarse en la selección, de acuerdo con la normativa de matrícula vigente.
\end{itemize}



